\documentclass[twocolumn,10pt,letterpaper]{article}

% -- Page geometry --
\usepackage[
  top=0.85in,
  bottom=0.75in,
  left=0.65in,
  right=0.65in,
  columnsep=0.25in
]{geometry}

% -- Beautiful typography: Century Schoolbook --
\usepackage[T1]{fontenc}
\usepackage{tgschola}              % TeX Gyre Schola (Century Schoolbook)
\usepackage[scaled=0.82]{beramono} % Bera Mono for code
\usepackage{microtype}             % Microtypographic refinements
\usepackage{setspace}              % Line spacing control
\setstretch{1.08}                  % Slightly open leading for readability

% -- Packages --
\usepackage{xurl}        % Allow URL breaks at any character (load before hyperref)
\usepackage{hyperref}
\usepackage{graphicx}
\usepackage{booktabs}    % Professional tables
\usepackage{tabularx}    % Width-constrained tables
\usepackage{enumitem}    % List customization
\usepackage{fancyhdr}
\usepackage{xcolor}
\usepackage{balance}
\usepackage{stfloats}    % Enable [h] placement for table* in two-column layout
\usepackage{titlesec}    % Section heading customization
\usepackage{listings}    % Code blocks with line wrapping

% -- Paragraph spacing (no indent, modest skip) --
\setlength{\parindent}{0pt}
\setlength{\parskip}{0.4em plus 0.1em minus 0.05em}
\setlength{\emergencystretch}{3em}  % Extra stretch to avoid overfull hboxes in narrow columns

% -- Code listings with line wrapping --
\lstset{
  basicstyle=\small\ttfamily,
  breaklines=true,
  breakatwhitespace=false,
  breakautoindent=true,
  postbreak=\mbox{\textcolor{accentrule}{$\hookrightarrow$}\space},
  columns=fullflexible,
  keepspaces=true,
  backgroundcolor=\color{codebg},
  frame=single,
  rulecolor=\color{codebg},
  framesep=3pt,
  aboveskip=0.5em,
  belowskip=0.5em,
  xleftmargin=4pt,
  xrightmargin=0pt,
}

% -- Colors --
\definecolor{linkblue}{HTML}{1D4ED8}
\definecolor{headergray}{HTML}{1F2937}
\definecolor{rulegray}{HTML}{D1D5DB}
\definecolor{titlebg}{HTML}{111827}
\definecolor{accentrule}{HTML}{6B7280}
\definecolor{brandred}{HTML}{8B2635}
\definecolor{codebg}{HTML}{F3F4F6}    % Very light gray for code backgrounds

% -- Hyperlinks --
\hypersetup{
    colorlinks=true,
    linkcolor=linkblue,
    urlcolor=linkblue,
    citecolor=linkblue,
    pdfstartview=FitH
}

% -- Section numbering: Roman > Alph > Arabic > alph, no decimals --
\renewcommand{\thesection}{\Roman{section}}
\renewcommand{\thesubsection}{\Alph{subsection}}
\renewcommand{\thesubsubsection}{\arabic{subsubsection}}
\renewcommand{\theparagraph}{\alph{paragraph}}
\makeatletter
\@addtoreset{subsection}{section}
\@addtoreset{subsubsection}{subsection}
\@addtoreset{paragraph}{subsubsection}
\makeatother

% -- Section headings --
\titleformat{\section}
  {\large\bfseries\color{titlebg}}
  {\thesection.}{0.5em}{}
  [\vspace{-0.3em}{\color{accentrule}\rule{\columnwidth}{0.5pt}}]
\titleformat{\subsection}
  {\normalsize\bfseries\color{headergray}}
  {\thesubsection.}{0.4em}{}
\titleformat{\subsubsection}
  {\small\bfseries\itshape\color{headergray}}
  {\thesubsubsection.}{0.4em}{}
\titleformat{\paragraph}[runin]
  {\small\bfseries\color{headergray}}
  {\theparagraph)}{0.4em}{}[.\hspace{0.4em}]
\titlespacing*{\section}{0pt}{1.0em}{0.35em}
\titlespacing*{\subsection}{0pt}{0.75em}{0.2em}
\titlespacing*{\subsubsection}{0pt}{0.55em}{0.15em}
\titlespacing*{\paragraph}{0pt}{0.4em}{0em}

% -- Lists: compact, elegant --
\setlist{nosep, leftmargin=1.1em, topsep=0.2em, itemsep=0.05em}
\setlist[itemize]{label={\small\textcolor{accentrule}{\textbullet}}}
\setlist[enumerate]{label={\small\textcolor{headergray}{\arabic*.}}}

% -- Tables: tighter, small text --
\renewcommand{\arraystretch}{1.15}

% -- Header/footer (pages 2+) --
\pagestyle{fancy}
\fancyhf{}
\renewcommand{\headrulewidth}{0.3pt}
\renewcommand{\headrule}{\hbox to\headwidth{\color{rulegray}\leaders\hrule height \headrulewidth\hfill}}
\fancyhead[L]{\footnotesize\color{headergray}\textit{Fork Progress Report}}
\renewcommand{\footrulewidth}{0.3pt}
\renewcommand{\footrule}{\hbox to\headwidth{\color{rulegray}\leaders\hrule height \footrulewidth\hfill}}
% Brand mark footer: CaseMirror wordmark left, page number right (both burgundy)
\fancyfoot[L]{\footnotesize\textbf{\color{headergray}CaseMirror}}
\fancyfoot[R]{\footnotesize\color{headergray}\thepage}

% -- Page 1 style: footer only, no header rule --
\fancypagestyle{firstpage}{%
  \fancyhf{}%
  \renewcommand{\headrulewidth}{0pt}%
  \renewcommand{\footrulewidth}{0.3pt}%
  \renewcommand{\footrule}{\hbox to\headwidth{\color{rulegray}\leaders\hrule height 0.3pt\hfill}}%
  \fancyfoot[L]{\footnotesize\textbf{\color{headergray}CaseMirror}}%
  \fancyfoot[R]{\footnotesize\color{headergray}\thepage}%
}

% -- Column rule --
\setlength{\columnseprule}{0.15pt}
\renewcommand{\columnseprulecolor}{\color{rulegray}}

% -- Title --
\title{\vspace{-1.8em}\LARGE\bfseries\color{titlebg} Fork Progress Report\vspace{-0.3em}}
\author{CaseMirror Research \\ \small\itshape\href{https://casemirror.ai}{\textcolor{headergray}{casemirror.ai}}}
\date{{\small\color{headergray} \today}}

\begin{document}

\maketitle
\thispagestyle{firstpage}
\vspace{-0.5em}


\textbf{Project:} \texttt{sia\_rust}  

\textbf{Repository:} \texttt{/\allowbreak home/\allowbreak m/\allowbreak sia\_rust}  

\textbf{Survey date:} 2026-06-06  

\textbf{Survey HEAD:} \texttt{da4191a} (\texttt{main})  


\subsection{Executive Summary}

Since the fork from the upstream SIA project, this repository has been transformed

from a Python-first self-improving-agent framework into an active Rust port with a

Python execution bridge, native Rust meta/feedback LLM runners behind an optional

\texttt{llm} feature, parity tests, a file-backed web dashboard, benchmark/eval harnesses,

and a substantial demo/reproducibility documentation package.


The fork-era work is concentrated in 48 commits on 2026-06-06 after the inferred

upstream baseline \texttt{4b4877f} (`chore: bump minor version as cli contracts have

changed (\#23)`). Across that range, the branch changes 139 files with 33,895

insertions and 333 deletions. The biggest areas of change are \texttt{src/\allowbreak }, \texttt{src/\allowbreak llm/\allowbreak },

\texttt{tests/\allowbreak }, \texttt{docs/\allowbreak }, \texttt{benchmarks/\allowbreak }, \texttt{evals/\allowbreak }, and the project issue-tracking state

under \texttt{.beads/\allowbreak }.


The project is in a strong Rust-port state for deterministic/offline surfaces:

default Rust tests pass, default and \texttt{llm} clippy pass, the standalone \texttt{evals/\allowbreak }

crate passes, and the Python/Rust differential parity script reports byte-identical

output on all checked surfaces. The main remaining risks are live-demo readiness,

live provider validation, hardening of the runtime security story, and replacing

some still-observational or reference-only "closed loop" pieces with acting,

production-quality behavior.


\subsection{Scope And Baseline}

There is no \texttt{upstream} remote configured in this checkout. The only configured

remote is:


\begin{itemize}
\setlength{\itemsep}{0pt}
\setlength{\parskip}{0pt}
  \item \texttt{origin}: \texttt{git@github.com:micahstubbs/\allowbreak sia\_rust.git}
\end{itemize}

Accordingly, this report treats the last pre-fork upstream-looking commit as an

inferred baseline rather than a remote-confirmed fork point:


\begin{itemize}
\setlength{\itemsep}{0pt}
\setlength{\parskip}{0pt}
  \item Upstream baseline used for this report: \texttt{4b4877f}
  \item Baseline subject: \texttt{chore: bump minor version as cli contracts have changed (#23)}
  \item Commits through that baseline: 11
  \item Fork-era commits after that baseline at survey time: 48
\end{itemize}

At survey time, \texttt{main} was ahead of \texttt{origin/\allowbreak main} by one commit:


\begin{itemize}
\setlength{\itemsep}{0pt}
\setlength{\parskip}{0pt}
  \item \texttt{da4191a} - \texttt{Sync beads tracker after upstream review}
\end{itemize}

During PDF generation, \texttt{origin/\allowbreak main} advanced by three additional remote commits

(\texttt{18e2f24}, \texttt{9442ced}, and \texttt{f1890b1}). This report evaluates the local checkout at

\texttt{da4191a} and does not include those later remote commits.


The report/PDF artifact commit generated from this document is not included in

those counts.


\subsection{Work Completed Since The Fork}

\subsubsection{Rust Port Foundation}

The largest milestone is \texttt{49415c2}, `Rust port of SIA framework (exhaustive parity

+ benchmarks + dsrs evals)`. It introduced the Rust crate, CLI, CI workflow,

benchmarks, tests, parity tooling, and the core ported modules.


Major Rust modules now mirror the Python package surface:


\begin{itemize}
\setlength{\itemsep}{0pt}
\setlength{\parskip}{0pt}
  \item Configuration, providers, profiles, API-key resolution, and bundled config files.
  \item Task layout and run-directory setup.
  \item Prompt building and context management.
  \item Orchestrator and generation loop scaffolding.
  \item Target-agent execution through a Python subprocess using the existing
\end{itemize}
  \texttt{evaluate.py} contract.

\begin{itemize}
\setlength{\itemsep}{0pt}
\setlength{\parskip}{0pt}
  \item Web data model and Axum server for \texttt{sia web}.
  \item CPython-compatible JSON formatting for byte-for-byte parity.
\end{itemize}

The public project framing was also rewritten. \texttt{README.md} now presents this as a

Rust port of SIA with a Python bridge rather than as the original Python package.

\texttt{docs/\allowbreak RUST\_PORT.md} documents the Python-to-Rust module map, build/test commands,

native-runner feature gate, parity surfaces, benchmarks, evals, and known testing

seams.


\subsubsection{Native Meta/Feedback LLM Runners}

The fork added a substantial native LLM layer under \texttt{src/\allowbreak llm/\allowbreak }, gated by the

optional \texttt{llm} Cargo feature. This includes:


\begin{itemize}
\setlength{\itemsep}{0pt}
\setlength{\parskip}{0pt}
  \item \texttt{AgentRunner} abstraction and trajectory context/outcome types.
  \item Anthropic Messages API transport and OpenAI-compatible chat transport types.
  \item Claude, OpenHands-style, and PydanticAI-style runner loops.
  \item Tool execution for read/write/edit/glob/bash.
  \item Trajectory logging middleware.
  \item Retry/backoff and transport decorators.
  \item Structured-output extraction/parity support.
  \item Provider/profile-to-client mapping.
  \item Per-run telemetry generation.
  \item Tavily web-search client primitive.
\end{itemize}

This closes a major boundary from the initial port: meta/feedback agents can now

use native Rust LLM clients when the crate is built with \texttt{-\allowbreak -\allowbreak features llm}, while

the target agent remains a Python subprocess so generated task agents preserve the

paper's task/evaluator contract.


\subsubsection{Closed-Loop Extensions: Scheduler, Weights, And Verifiers}

The fork adds first-pass closed-loop machinery beyond the baseline SIA harness

update flow:


\begin{itemize}
\setlength{\itemsep}{0pt}
\setlength{\parskip}{0pt}
  \item Adaptive scheduler logic in \texttt{src/\allowbreak scheduler.rs}.
  \item Scheduler decision recording and dashboard surfacing through \texttt{src/\allowbreak closed\_loop.rs}
\end{itemize}
  and web data endpoints.

\begin{itemize}
\setlength{\itemsep}{0pt}
\setlength{\parskip}{0pt}
  \item A native \texttt{WeightUpdater} abstraction and CPU reference LoRA-style updater in
\end{itemize}
  \texttt{src/\allowbreak weights.rs}.

\begin{itemize}
\setlength{\itemsep}{0pt}
\setlength{\parskip}{0pt}
  \item Reusable verifier traits and task robustness helpers in \texttt{src/\allowbreak verifier.rs}.
\end{itemize}

These are meaningful research/product extensions, but the issue tracker correctly

records that some are still partial. In particular, the current scheduler path is

observational in key places, and the CPU weight-update path is a reference toy

rather than real model training or an adapter that is loaded by the target agent.


\subsubsection{SIA Studio And Web/Demo Surface}

The fork upgraded the web surface into a more complete "SIA Studio" dashboard:


\begin{itemize}
\setlength{\itemsep}{0pt}
\setlength{\parskip}{0pt}
  \item Telemetry and metrics charts.
  \item Dark-mode/polished UI work.
  \item Run-level scheduler decision timeline.
  \item API endpoints for telemetry, metrics, scheduler decisions, and weight updates.
\end{itemize}

The dashboard is functional for file-backed run visualization, but current open

issues identify important demo gaps: auto-refresh is missing, port-bind failures

can be hidden, and run IDs can collide during rehearsal or stage demos.


\subsubsection{Security And Sandboxing}

The fork added a formal security posture and first hardening primitives:


\begin{itemize}
\setlength{\itemsep}{0pt}
\setlength{\parskip}{0pt}
  \item \texttt{SECURITY.md} threat model.
  \item \texttt{src/\allowbreak sandbox.rs} capability allow-list abstraction.
  \item Tests for read/write/bash permissions, path containment, and size limits.
  \item Native tool executors wired to capability checks.
  \item Docker sandbox command generation for Python target-agent execution.
\end{itemize}

The current security story is honest but not finished. The default native-runner

capability profile still allows broad Bash behavior, and the Docker sandbox path is

not yet a drop-in safe path for live LLM-calling target agents because network,

dependencies, and credentials require an explicit policy.


\subsubsection{Providers, Credentials, And Reproducibility}

The fork added or expanded provider support and documentation:


\begin{itemize}
\setlength{\itemsep}{0pt}
\setlength{\parskip}{0pt}
  \item Bundled Nebius profiles and model-slug verification tooling.
  \item Bundled Gemini target profile and tests.
  \item Bundled Tinker provider plus \texttt{gptoss-\allowbreak tinker-\allowbreak target} and
\end{itemize}
  \texttt{qwen3-\allowbreak tinker-\allowbreak target} profiles.

\begin{itemize}
\setlength{\itemsep}{0pt}
\setlength{\parskip}{0pt}
  \item \texttt{.env} loading at startup, with real environment variables taking precedence.
  \item Per-provider credentials documentation in \texttt{docs/\allowbreak CREDENTIALS.md}.
  \item Nebius quickstart and model verification docs.
  \item Reproducibility standards and run-artifact checklist.
  \item Hackathon demo script, slide outline, run-of-show, and judging narrative.
  \item Preprint-style draft in \texttt{docs/\allowbreak paper/\allowbreak sia\_rust\_preprint.md}.
\end{itemize}

This gives the project a strong external-facing package, but live provider paths

remain mostly guarded behind ignored tests that require real keys and network

access.


\subsubsection{Benchmarks And Evals}

The fork added multiple validation/evaluation assets:


\begin{itemize}
\setlength{\itemsep}{0pt}
\setlength{\parskip}{0pt}
  \item Criterion benchmark scaffold under \texttt{benches/\allowbreak core.rs}.
  \item Python-vs-Rust benchmark scripts under \texttt{benchmarks/\allowbreak }.
  \item \texttt{benchmarks/\allowbreak REPORT.md}, documenting the deterministic core as about 5.8x faster
\end{itemize}
  by geometric mean in the benchmark report.

\begin{itemize}
\setlength{\itemsep}{0pt}
\setlength{\parskip}{0pt}
  \item Standalone \texttt{evals/\allowbreak } crate using DSRs/dspy-rs for a GPQA-style multiple-choice
\end{itemize}
  evaluation harness with offline mock tests.

\begin{itemize}
\setlength{\itemsep}{0pt}
\setlength{\parskip}{0pt}
  \item Differential parity script at \texttt{scripts/\allowbreak parity\_check.py}.
\end{itemize}

\subsubsection{Project Operations And Issue Tracking}

The latest local commits added project operations scaffolding:


\begin{itemize}
\setlength{\itemsep}{0pt}
\setlength{\parskip}{0pt}
  \item \texttt{.beads/\allowbreak } project configuration and issue export.
  \item \texttt{CLAUDE.md} local agent guidance.
  \item \texttt{package.json}.
  \item \texttt{scripts/\allowbreak seed\_beads\_from\_gh.sh}.
  \item A session summary under \texttt{docs/\allowbreak session-\allowbreak summaries/\allowbreak }.
  \item Mirroring of 18 open GitHub issues into Beads.
  \item Subsequent Beads synchronization after upstream review.
  \item Expanded \texttt{CLAUDE.md} local agent instructions covering build, test, parity, and
\end{itemize}
  architecture guidance.

\begin{itemize}
\setlength{\itemsep}{0pt}
\setlength{\parskip}{0pt}
  \item Python packaging metadata guard for the OpenHands optional extra, constraining
\end{itemize}
  it to supported Python versions.


At survey time, \texttt{br list} showed 18 open issues. \texttt{br ready -\allowbreak -\allowbreak limit 20} showed 17

ready issues with no blockers. The P1 ready work is heavily demo/safety oriented:


\begin{itemize}
\setlength{\itemsep}{0pt}
\setlength{\parskip}{0pt}
  \item Harden native-runner capability profile for demo safety.
  \item Make Docker sandbox usable or clearly scoped for LLM-calling target agents.
  \item Add a stage-safe live demo command.
  \item Make SIA Studio auto-refresh during live runs.
\end{itemize}

\subsection{Current Verification Results}

Commands run during this survey:


\begin{itemize}
\setlength{\itemsep}{0pt}
\setlength{\parskip}{0pt}
  \item Default Rust suite: passed with \texttt{cargo test}; the library crate reported 154
\end{itemize}
  passed tests, and the integration/doc tests also passed.

\begin{itemize}
\setlength{\itemsep}{0pt}
\setlength{\parskip}{0pt}
  \item Parallel \texttt{llm} suite: failed under default parallel execution; the library
\end{itemize}
  crate reported 272 passed, 4 failed, and 6 ignored tests before aborting. The

  failures were \texttt{llm::provider\_mapping} tests racing on process environment

  variables.

\begin{itemize}
\setlength{\itemsep}{0pt}
\setlength{\parskip}{0pt}
  \item Serialized \texttt{llm} suite: passed with \texttt{-\allowbreak -\allowbreak test-\allowbreak threads=1}; the library crate
\end{itemize}
  reported 276 passed and 6 ignored tests, and the \texttt{llm} integration/doc tests

  also passed.

\begin{itemize}
\setlength{\itemsep}{0pt}
\setlength{\parskip}{0pt}
  \item Standalone eval crate: passed with 4 offline tests.
  \item Formatting and lints: \texttt{cargo fmt -\allowbreak -\allowbreak check}, default clippy, and \texttt{llm} clippy all
\end{itemize}
  passed with warnings denied.

\begin{itemize}
\setlength{\itemsep}{0pt}
\setlength{\parskip}{0pt}
  \item Differential parity: passed; the script reported `PARITY OK: all surfaces
\end{itemize}
  byte-identical`.

\begin{itemize}
\setlength{\itemsep}{0pt}
\setlength{\parskip}{0pt}
  \item Python pytest suite: not run because \texttt{pytest} is not installed in this
\end{itemize}
  environment.

\begin{itemize}
\setlength{\itemsep}{0pt}
\setlength{\parskip}{0pt}
  \item Packaging metadata unittest: passed.
\end{itemize}

The normal parallel \texttt{llm} test failure is worth treating as a real test-isolation

issue even though serialized execution passes. The failing tests mutate global

environment variables using a helper that snapshots/restores state, but other

parallel tests also read or mutate those variables. A mutex or serial-test pattern

would make this robust under default \texttt{cargo test -\allowbreak -\allowbreak features llm}.


\subsection{Current Worktree State}

At survey time, before updating this document and generating the PDF artifacts,

the worktree was clean. \texttt{main} was one commit ahead of \texttt{origin/\allowbreak main}, with the

only ahead commit being \texttt{da4191a} (\texttt{Sync beads tracker after upstream review}).

During report/PDF generation, an unstaged \texttt{.beads/\allowbreak issues.jsonl} update appeared;

it is intentionally outside the report artifact commit.


\subsection{Main Remaining Risks}

The current backlog and verification results point to these high-value next steps:


\begin{enumerate}
\setlength{\itemsep}{0pt}
\setlength{\parskip}{0pt}
  \item Fix the \texttt{llm} feature test race so \texttt{cargo test -\allowbreak -\allowbreak features llm} passes under
\end{enumerate}
   default parallel test execution.

\begin{enumerate}
\setlength{\itemsep}{0pt}
\setlength{\parskip}{0pt}
  \item Run live provider validation with real keys: ignored Anthropic, OpenHands,
\end{enumerate}
   PydanticAI, Gemini, Nebius, structured-output, and Tavily live tests.

\begin{enumerate}
\setlength{\itemsep}{0pt}
\setlength{\parskip}{0pt}
  \item Run a complete live \texttt{sia run -\allowbreak -\allowbreak features llm} self-improvement loop and capture
\end{enumerate}
   the artifacts.

\begin{enumerate}
\setlength{\itemsep}{0pt}
\setlength{\parskip}{0pt}
  \item Make the scheduler decision actually drive the generation loop rather than only
\end{enumerate}
   producing artifacts/logging.

\begin{enumerate}
\setlength{\itemsep}{0pt}
\setlength{\parskip}{0pt}
  \item Replace or supplement the CPU reference weight updater with a real model update
\end{enumerate}
   path, or clearly scope it as an offline demonstrator.

\begin{enumerate}
\setlength{\itemsep}{0pt}
\setlength{\parskip}{0pt}
  \item Harden native-runner tool permissions for demo mode.
  \item Clarify or implement Docker sandbox support for live provider-calling target
\end{enumerate}
   agents.

\begin{enumerate}
\setlength{\itemsep}{0pt}
\setlength{\parskip}{0pt}
  \item Fix demo fragility: SIA Studio auto-refresh, dashboard port handling,
\end{enumerate}
   collision-proof run IDs, and a fast stage-safe task.

\begin{enumerate}
\setlength{\itemsep}{0pt}
\setlength{\parskip}{0pt}
  \item Resolve documentation/package identity issues that still point to the original
\end{enumerate}
   Python repository.

\begin{enumerate}
\setlength{\itemsep}{0pt}
\setlength{\parskip}{0pt}
  \item Install/use the Python test environment so the full Python pytest suite can be
\end{enumerate}
    included in regular verification.


\subsection{Overall Assessment}

The fork has already accomplished the hard architectural work: the deterministic

SIA loop is represented in Rust, parity is actively tested, native meta/feedback

LLM runners exist, telemetry and dashboard surfaces are in place, and the project

has a strong documentation and demo package.


The project is not yet "live-demo hardened." Its deterministic and offline

verification story is strong, but the live-provider path, security defaults,

dashboard freshness, and closed-loop actuation still need focused work before the

repository can credibly claim a fully production-ready self-improving loop.



\balance
\end{document}
